% GENERATED FILE. Edit canonical lessons, then run npm run generate:books.
% canonical-source-hash: fnv1a64:bb322c025ce7c59e
% canonical-lessons: LA-C181-gratuitus, LA-C181-recens, LA-C181-crudus, LA-C181-coctus, LA-C181-longinquus

\chapter{Free of charge, Fresh, Raw, Cooked, Distant}
\label{ch:la-free-of-charge-fresh-raw-cooked-distant}
\hlchaptermodality{181}

\begin{chapteropening}
\textbf{By the end of this chapter:} \emph{I can say free of charge, fresh, raw, cooked and distant.}

Complete the last lesson of chapter 181: I can say free of charge, fresh, raw, cooked and distant.
\end{chapteropening}

\section[grātuītus, grātuīta, grātuītum]{grātuītus, grātuīta, grātuītum — free of charge}

\label{lesson:LA-C181-gratuitus}



\begin{quote}
Before the new one: say the Latin for cheerful, then the Latin for shy.
\end{quote}

\subsection*{You'll want to know: grātuītus, grātuīta, grātuītum}
\textbf{grātuītus, grātuīta, grātuītum} — \textquotedblleft{}free of charge\textquotedblright{}.

\begin{cousinweb}
\textbf{grātuītus} is built on \emph{grātia}, \textquotedblleft{}favour\textquotedblright{} (Lewis \& Short): done as a favour, without pay. English \emph{gratuitous} and \emph{gratis} belong to the same family.
\end{cousinweb}

\begin{grammarlens}[title={What you've built}]
One more word for this chapter.
\end{grammarlens}

\subsection*{Your turn}
\begin{itemize}
\raggedright
  \item \emph{Say it:} \emph{grātuītus, grātuīta, grātuītum}
  \item \emph{Say it:} \emph{grātuītus, grātuīta, grātuītum}, once more
  \item \emph{Say it:} say \emph{grātuītus, grātuīta, grātuītum}
  \item \emph{Recall:} say the Latin for cheerful, then the Latin for shy, then say \emph{grātuītus, grātuīta, grātuītum} again
\end{itemize}

\begin{tcolorbox}[breakable,colback=teal!4,colframe=teal!35!black,title={Before you move on}]
What does grātuītus, grātuīta, grātuītum mean? (\textquotedblleft{}Free of charge\textquotedblright{}.) Say it once more.
\end{tcolorbox}
\section[recēns, recentis]{recēns, recentis — fresh, recent}

\label{lesson:LA-C181-recens}



\begin{quote}
Before the new one: say the Latin for shy, then the Latin for free of charge.
\end{quote}

\subsection*{You'll want to know: recēns, recentis}
\textbf{recēns, recentis} — \textquotedblleft{}fresh, recent\textquotedblright{}.

\begin{cousinweb}
\textbf{recēns} has no sure origin; Greek \emph{kainos}, \textquotedblleft{}new\textquotedblright{}, is often compared (Lewis \& Short), but the details are uncertain.
\end{cousinweb}

\begin{grammarlens}[title={What you've built}]
One more word for this chapter.
\end{grammarlens}

\subsection*{Your turn}
\begin{itemize}
\raggedright
  \item \emph{Say it:} \emph{recēns, recentis}
  \item \emph{Say it:} \emph{recēns, recentis}, once more
  \item \emph{Say it:} say \emph{recēns, recentis}
  \item \emph{Recall:} say the Latin for shy, then the Latin for free of charge, then say \emph{recēns, recentis} again
\end{itemize}

\begin{tcolorbox}[breakable,colback=teal!4,colframe=teal!35!black,title={Before you move on}]
What does recēns, recentis mean? (\textquotedblleft{}Fresh, recent\textquotedblright{}.) Say it once more.
\end{tcolorbox}
\section[crūdus, crūda, crūdum]{crūdus, crūda, crūdum — raw}

\label{lesson:LA-C181-crudus}



\begin{quote}
Before the new one: say the Latin for free of charge, then the Latin for fresh.
\end{quote}

\subsection*{You'll want to know: crūdus, crūda, crūdum}
\textbf{crūdus, crūda, crūdum} — \textquotedblleft{}raw\textquotedblright{}.

\begin{cousinweb}
\textbf{crūdus} is inherited, usually traced to Proto-Indo-European \emph{*krewh2-}, \textquotedblleft{}raw flesh, blood\textquotedblright{} (EtymDB 2.1); \emph{cruor}, \textquotedblleft{}blood\textquotedblright{}, is its cousin in Latin (Lewis \& Short), and English \emph{raw} is a cousin too.
\end{cousinweb}

\begin{grammarlens}[title={What you've built}]
One more word for this chapter.
\end{grammarlens}

\subsection*{Your turn}
\begin{itemize}
\raggedright
  \item \emph{Say it:} \emph{crūdus, crūda, crūdum}
  \item \emph{Say it:} \emph{crūdus, crūda, crūdum}, once more
  \item \emph{Say it:} say \emph{crūdus, crūda, crūdum}
  \item \emph{Recall:} say the Latin for free of charge, then the Latin for fresh, then say \emph{crūdus, crūda, crūdum} again
\end{itemize}

\begin{tcolorbox}[breakable,colback=teal!4,colframe=teal!35!black,title={Before you move on}]
What does crūdus, crūda, crūdum mean? (\textquotedblleft{}Raw\textquotedblright{}.) Say it once more.
\end{tcolorbox}
\section[coctus, cocta, coctum]{coctus, cocta, coctum — cooked}

\label{lesson:LA-C181-coctus}



\begin{quote}
Before the new one: say the Latin for fresh, then the Latin for raw.
\end{quote}

\subsection*{You'll want to know: coctus, cocta, coctum}
\textbf{coctus, cocta, coctum} — \textquotedblleft{}cooked\textquotedblright{}.

\begin{cousinweb}
\textbf{coctus} is the past participle of \emph{coquō}, \textquotedblleft{}to cook\textquotedblright{} (Lewis \& Short): cooked.
\end{cousinweb}

\begin{grammarlens}[title={What you've built}]
One more word for this chapter.
\end{grammarlens}

\subsection*{Your turn}
\begin{itemize}
\raggedright
  \item \emph{Say it:} \emph{coctus, cocta, coctum}
  \item \emph{Say it:} \emph{coctus, cocta, coctum}, once more
  \item \emph{Say it:} say \emph{coctus, cocta, coctum}
  \item \emph{Recall:} say the Latin for fresh, then the Latin for raw, then say \emph{coctus, cocta, coctum} again
\end{itemize}

\begin{tcolorbox}[breakable,colback=teal!4,colframe=teal!35!black,title={Before you move on}]
What does coctus, cocta, coctum mean? (\textquotedblleft{}Cooked\textquotedblright{}.) Say it once more.
\end{tcolorbox}
\section[longinquus, longinqua, longinquum]{longinquus, longinqua, longinquum — distant}

\label{lesson:LA-C181-longinquus}



\begin{quote}
Before the new one: say the Latin for raw, then the Latin for cooked.
\end{quote}

\subsection*{You'll want to know: longinquus, longinqua, longinquum}
\textbf{longinquus, longinqua, longinquum} — \textquotedblleft{}distant\textquotedblright{}.

\begin{cousinweb}
\textbf{longinquus} is built on \emph{longus}, \textquotedblleft{}long\textquotedblright{} (Lewis \& Short), with an ending also found in \emph{propinquus}, \textquotedblleft{}near\textquotedblright{}: far off.
\end{cousinweb}

\begin{grammarlens}[title={What you've built}]
One more word for this chapter.
\end{grammarlens}

\subsection*{Your turn}
\begin{itemize}
\raggedright
  \item \emph{Say it:} \emph{longinquus, longinqua, longinquum}
  \item \emph{Say it:} \emph{longinquus, longinqua, longinquum}, once more
  \item \emph{Say it:} say \emph{longinquus, longinqua, longinquum}
  \item \emph{Recall:} say the Latin for raw, then the Latin for cooked, then say \emph{longinquus, longinqua, longinquum} again
\end{itemize}

\begin{tcolorbox}[breakable,colback=teal!4,colframe=teal!35!black,title={Before you move on}]
What does longinquus, longinqua, longinquum mean? (\textquotedblleft{}Distant\textquotedblright{}.) Say it once more.
\end{tcolorbox}
